\documentclass{article}
\usepackage{graphicx} % Required for inserting images
\usepackage{amsmath}
\title{Lesson 24 Homework - Proof Writing}
\author{William Wang}
\date{February 2026}

\begin{document}

\maketitle

\section{}


Let \(P=(x,y)\). We first want to reflect over the y axis. Reflection over the \(y\)-axis becomes
\[(x,y)\longrightarrow(-x,y).\]
Then we want to rotate by 60 degrees counterclockwise. The rotation matrix for \(60^\circ\) is
\[
A =
\begin{pmatrix}
\cos 60^\circ & -\sin 60^\circ \\
\sin 60^\circ & \cos 60^\circ
\end{pmatrix}=
\begin{pmatrix}\frac12 & -\frac{\sqrt3}{2} \\\frac{\sqrt3}{2} & \frac12\end{pmatrix}\]

Applying this to \((-x,y)\):

\[
\begin{aligned}
x' &= (-x)\frac12 - y\frac{\sqrt3}{2}
= -\frac{x}{2} - \frac{\sqrt3}{2}y, \\
y' &= (-x)\frac{\sqrt3}{2} + y\frac12
= -\frac{\sqrt3}{2}x + \frac12 y.
\end{aligned}
\]

We are given
\[
(x',y')=(-2-\sqrt3,\;1-2\sqrt3).
\]

So the system is
\[
-\frac{x}{2}-\frac{\sqrt3}{2}y=-2-\sqrt3,
\]
\[
-\frac{\sqrt3}{2}x+\frac12 y=1-2\sqrt3.
\]

\[
-x-\sqrt3 y=-4-2\sqrt3,
\]
\[
-\sqrt3 x+y=2-4\sqrt3.
\]
\[
y=2-4\sqrt3+\sqrt3 x.
\]

Substitute into the first:

\[
-x-\sqrt3(2-4\sqrt3+\sqrt3 x)=-4-2\sqrt3.
\]

\[
-4x+12=-4.
\]

\[
-4x=-16 \quad\Rightarrow\quad x=4.
\]

\[
y=2-4\sqrt3+4\sqrt3=2.
\]

\[
P=(4,2)
\]

\section{}
If we think in terms of geometric rotation, then we can view \(A^2\) as a rotation by \(2\theta\) and \(A^3\) as a rotation of \(3\theta\). Which means we would have a rotation of \(4\theta\) the same as a rotation of \(\theta\). 
\[4\theta \equiv \theta \pmod{2\theta}\]

\[3\theta \equiv 0 \pmod{2\theta}\]
\[\theta = 2\pi/3\]
Hence the the result is the matrix for the transformation of rotating by \(2\pi/3\). 
\[\begin{pmatrix}
    -1/2&-\sqrt{3}/2 \\ \sqrt{3}/2&-1/2
\end{pmatrix}\]

\section{}
We first want to find a position vector that is on the line \(y = 2x\). \(P = \langle 1,2 \rangle\). We want to find the projection of the vector towards the standard basis vectors. 
\[\text{proj}_{\langle  1,2 \rangle}^{\langle 1,0 \rangle} = \frac{1}{5}\langle 1,2 \rangle = \langle \frac{1}{5}, \frac{2}{5} \rangle\]
Then to get the normal vector, we want to subtract this from the standard basis vector. 
\[\vec{n_1} = \langle 1,0\rangle - \langle 1/5, 2/5 \rangle = \langle  4/5, -2/5\rangle\]
Then for the other standard basis vector
\[\text{proj}_{\langle 1,2 \rangle}^{\langle 0,1 \rangle} = \frac{2}{5}\langle 1,2 \rangle = \langle \frac{2}{5}, \frac{4}{5} \rangle \]
\[\vec{n_2} = \langle 0,1 \rangle - \rangle 2/5, 4/5\langle = \langle -2/5, 1/5 \rangle  \]
Then we have the matrix 
\[A = \begin{pmatrix}
    4/5&-2/5 \\ -2/5&1/5
\end{pmatrix}\]
which means the transformation is 
\[T\begin{pmatrix}
    x \\ y
\end{pmatrix} = \begin{pmatrix}
    4/5&-2/5 \\ -2/5&1/5
\end{pmatrix}\begin{pmatrix}
    x \\ y
\end{pmatrix}\]
The transformation of \((4,-6)\) will be 
\[\begin{pmatrix}
    4/5&-2/5 \\ -2/5&1/5
\end{pmatrix}\begin{pmatrix}
    4 \\ -6
\end{pmatrix} = \begin{pmatrix}
    28/5 \\ -14/5
\end{pmatrix}\]
Hence 
\[(28/5,-14/5)\]
\end{document}
